\section{Nonvanishing when the minimal intersection lattice has rank at least two}
\label{sec:nonrank}

Let $X_0$ be a minimal surface of general type with $p_g=q=0$, and write
$K_0=K_{X_0}$ and $k_0=K_0^2$.  By Proposition~\ref{prop:reduction},
$K_0$ is nef and big and $1\leq k_0\leq9$.  This Section treats
$k_0\leq8$.

\begin{proposition}\label{prop:nonrank}
Suppose $1\leq k_0\leq8$.  Blow up
\[
 e=\max\{k_0-5,0\}
\]
distinct points of $X_0$, and call the resulting surface $X$.
There are choices of $L_1,L_2,D,n,p$ on $X$ for which
Theorem~\ref{thm:chow-relation} gives a relation with $c>0$.
Consequently the degree map on $\CH_0(X_0)_{\mathbb Q}$ is an isomorphism.
\end{proposition}

The proof pairs the Chow relation with $\alpha\boxtimes\alpha$, for a
class $\alpha\in H^2(X,\mathbb R)$ with $\alpha K_X=0$ and
$\alpha^2<0$.  We will take $L_1,L_2,D$ to be integral multiples of
$K_X$.  A fixed contribution indexed by $T+H=D$ with $\alpha H=0$
then pairs to $c_*\alpha^2$, since its external products pair to zero.
We arrange that every remaining contribution has zero point lengths,
hence $c_*=0$, and has positive pairing.  At least one such contribution
will occur, so the vanishing of the full pairing will force $c>0$.
The lattice calculation supplies $\alpha$; the subsequent numerical
bounds isolate the contributions with zero point lengths.

\subsection{A characteristic-vector calculation}

The exponential sequence and $H^1(\mathcal O_{X_0})=
H^2(\mathcal O_{X_0})=0$ identify $\Pic(X_0)$ with
$H^2(X_0,\mathbb Z)$.  In particular every integral cohomology class is
the first Chern class of a line bundle.  The free quotient
\[
 \Lambda=H^2(X_0,\mathbb Z)/\mathrm{torsion}
\]
is unimodular by Poincar\'e duality.  Its signature is $(1,b)$ with
$b=9-k_0$, by Hodge index and $b_2(X_0)=10-k_0$.
Riemann--Roch gives
\[
 x^2\equiv K_0x\pmod2\qquad(x\in\Lambda),
\]
so $K_0$ is a characteristic vector.

We use the classification of indefinite unimodular integral forms
\cite[Chapter~V, \S2.1, Corollary~1 to Theorem~2, and \S2.2,
Theorems~4--6]{SerreArithmetic}:
an odd form of signature $(1,b)$ is
$I_{1,b}=\langle1\rangle\oplus\langle-1\rangle^{\oplus b}$;
the signature difference of an even unimodular form is divisible by
eight.  Since $1\leq b\leq8$, the even case has $b=1$ and is the
hyperbolic plane.

\begin{lemma}\label{nr:lattice-class}
If $\Lambda$ is odd, there is a line bundle $h_0$ on $X_0$ with
\[
 h_0^2=-1,\qquad K_0h_0=1.
\]
If $\Lambda$ is even, then $k_0=8$ and there is a line bundle $h_0$ with
\[
 h_0^2=0,\qquad K_0h_0=2.
\]
\end{lemma}

\begin{proof}
In the odd case choose diagonal coordinates and write
$K_0=(u;v_1,\ldots,v_b)$.  The characteristic property says that all
these integers are odd.  Coordinate sign changes and permutations allow
us to assume
\[
 u>0,\qquad v_1\geq\cdots\geq v_b>0,
 \qquad u^2=9-b+\sum_{i=1}^b v_i^2.
\]
If $v_i=1$, the class $-e_i$, for the corresponding negative basis
vector $e_i$, has square $-1$ and pairing $1$ with $K_0$.
We show by descent on $u$ that this
case can always be reached by integral lattice isometries.

For $b=1$, the equation
$(u-v_1)(u+v_1)=8$ forces $(u,v_1)=(3,1)$: both factors are positive
even integers and the second is strictly larger than the first.
Suppose next that $b\geq3$ and that every $v_i\geq3$.  Put
$\ell=v_1+v_2+v_3$.  The ordering of the $v_i$ gives
\begin{align*}
 \ell^2-u^2
 &=2(v_1v_2+v_1v_3+v_2v_3)
       -\sum_{i=4}^b v_i^2-(9-b)\\
 &\geq (9-b)(v_3^2-1)>0.
\end{align*}
On the other hand,
$\ell^2\leq3(v_1^2+v_2^2+v_3^2)<3u^2<4u^2$.
The vector $z=(1;1,1,1,0,\ldots,0)$ has square $-2$.
Its integral reflection $x\mapsto x+(xz)z$ changes the positive
coordinate of $K_0$ to
\[
 u'=2u-\ell,\qquad 0<u'<u.
\]

For $b=2$ and $v_1,v_2\geq3$, put $\ell=v_1+v_2$.  Here
\[
 \ell^2-u^2=2v_1v_2-7>0,
 \qquad \ell^2\leq2(u^2-7)<2u^2<\frac94u^2.
\]
The vector $z=(1;1,1)$ has square $-1$, so its reflection is the
integral isometry $x\mapsto x+2(xz)z$.  It changes the positive
coordinate to
\[
 u'=3u-2\ell,\qquad 0<u'<u.
\]
In either case the reflected vector remains characteristic and of square
$9-b$.  Rechoose the signs and order of its negative coordinates and
repeat.  The positive integer $u$ strictly decreases, so the process
must reach a coordinate $v_i=1$.  Pulling the corresponding vector back
through the isometries gives the asserted class, which lifts to a line
bundle.

In the even case choose an isotropic basis $f,g$ with $fg=1$.
Writing $K_0=af+bg$, the characteristic property makes $a,b$ even,
and $2ab=K_0^2=8$.  After changing both basis signs if necessary,
$a=b=2$.  Either $f$ or $g$ gives $h_0$.
\end{proof}

Let $\pi:X\to X_0$ be the blowup in Proposition~\ref{prop:nonrank},
with exceptional curves $E_1,\ldots,E_e$.  We usually suppress $\pi^*$
on line and cohomology classes.  Thus
\[
 K=K_X=\pi^*K_0+\sum_{j=1}^e E_j,
 \qquad k=K^2=k_0-e=\min\{k_0,5\}>0.
\]
In the odd case put $h'=\pi^*h_0$.  In the even case $k_0=8$ and
$e=3$, so we may put $h'=\pi^*h_0+E_1$.  Both choices satisfy
\begin{equation}\label{nr:h-prime}
 (h')^2=-1,\qquad Kh'=1.
\end{equation}
The symbol $h'$ denotes a line bundle and is distinct from the insertion
class $h$ on the Quot scheme.

Set $a=K_0h_0$, so $a=1$ in the odd case and $a=2$ in the even case,
and define
\begin{equation}\label{nr:alpha}
 \alpha=\pi^*\!\left(h_0-\frac{a}{k_0}K_0\right)
       \in H^2(X,\mathbb R).
\end{equation}
Pullback classes are orthogonal to the exceptional curves, whence
\begin{equation}\label{nr:alpha-identities}
 \alpha K=0,\qquad
 \alpha h'=\alpha^2=h_0^2-\frac{a^2}{k_0}<0.
\end{equation}

\subsection{Positive canonical pairings on the blowup}

The following bound will apply to each of the two effective divisor
classes in a fixed component.

\begin{lemma}\label{nr:positive-pairing}
Let $P$ be an effective line bundle on $X$ such that $m_P\geq0$ and
$\alpha P\ne0$.  Then $KP>k$.
\end{lemma}

\begin{proof}
Write
\[
 P=\pi^*P_0+\sum_{j=1}^e a_jE_j,
 \qquad r_0=K_0P_0,
\]
where the $a_j$ are integers.  Pushing forward an effective divisor
representing $P$ shows that $P_0$ is effective; consequently
$r_0\geq0$ by nefness of $K_0$.
The inequality $2m_P=P^2-KP\geq0$ gives
\begin{equation}\label{nr:exceptional-bound}
 0\leq\sum_{j=1}^e a_j(a_j-1)
 \leq P_0^2-r_0
 <\frac{r_0^2}{k_0}-r_0.
\end{equation}
The first inequality holds for all integral $a_j$, including negative
ones.  The last is strict Hodge index: $\alpha P\ne0$ implies that
$P_0$ is not numerically proportional to $K_0$.  Together with
$r_0\geq0$, this forces $r_0>k_0$.

We next prove $KP>0$.  If $e=0$, this is immediate.  Otherwise put
$A_0=\sum_j a_j$ and suppose $KP=r_0-A_0\leq0$.  By Cauchy--Schwarz,
\[
 \sum_j a_j(a_j-1)
 \geq\frac{A_0^2}{e}-A_0
 \geq\frac{r_0^2}{e}-r_0
 >\frac{r_0^2}{k_0}-r_0.
\]
Indeed $A_0\geq r_0>k_0>e$, and $u\mapsto u^2/e-u$ is increasing
on $[r_0,\infty)$.  This contradicts
\eqref{nr:exceptional-bound}.  Thus $KP>0$.

Since $K^2=k>0$, the intersection form is negative definite on
$K^\perp$.  Also $\alpha P\ne0$ and $\alpha K=0$ show that $P$
is not numerically proportional to $K$.  Therefore
\[
 0<KP\leq P^2<\frac{(KP)^2}{k},
\]
where the middle inequality is $m_P\geq0$.  Dividing by $KP>0$
proves $KP>k$.  This application of Hodge index requires the positive
square of $K$; all effectivity bounds above used the nef class $\pi^*K_0$.
\end{proof}

\subsection{All contributions with nonzero divisor pairing}

We choose the Quot data so that every fixed contribution detected by
the divisor pairing with \(\alpha\) has zero point lengths. This will
make those contributions explicit, while the orthogonal terms account
for the diagonal coefficient. Set
\begin{equation}\label{nr:data}
 \begin{array}{c|cc}
  &d&p\\\hline
  k=1&4&0\\
  2\leq k\leq5&3&k-2
 \end{array}
 \qquad
 D=dK,\quad L_1=0,\quad L_2=K,\quad n=d(k+1).
\end{equation}
Here $A=-K$, so \eqref{eq:vdim} gives $v=2d$.
For $k=1$ we have $v=8>2=p+2$; for $2\leq k\leq5$ we have
$v=6>k=p+2$.  Thus Theorem~\ref{thm:chow-relation} applies.
This degree inequality is the reason for using blowups to reduce $k$
to at most five.

By Proposition~\ref{prop:fixed-pg0}, only effective line decompositions
$T+H=D$ with $m_T,m_H\geq0$ need be retained.  Write
\[
 r=KH,\qquad s=H^2,\qquad N=n_1+n_2.
\]
For these data the length constraint \eqref{eq:length} is
\begin{equation}\label{nr:length-specialized}
 N=d(k+1)-(d+1)r+s\geq0.
\end{equation}

\begin{lemma}\label{nr:surviving-numerics}
Every such fixed contribution with $\alpha H\ne0$ satisfies
\begin{equation}\label{nr:survivors}
 KH=H^2=k+1,\qquad
 n_1=n_2=0,\qquad m_H=0,\qquad m_T=p.
\end{equation}
\end{lemma}

\begin{proof}
Since $\alpha T=-\alpha H\ne0$, Lemma~\ref{nr:positive-pairing}
applies to both $H$ and $T$.  It gives
\[
 r>k,\qquad dk-r>k,\qquad r\leq s<\frac{r^2}{k}.
\]
Furthermore $s-r$ is even by Riemann--Roch.

If $k=1$, then $d=4$ and the first two bounds force $r=2$.
The inequalities $2\leq s<4$ and $s\equiv2\pmod2$ force $s=2$.
Equation~\eqref{nr:length-specialized} gives $N=0$.

For $2\leq k\leq5$, write $r=k+j$.  Since $d=3$, the two strict
canonical bounds give $1\leq j\leq k-1$.  Use the parity condition
by writing $s=r+2m$, where $m=m_H$ is a nonnegative integer.
The strict Hodge bound gives
\begin{equation}\label{nr:length-bound}
 2m<\frac{r^2}{k}-r=j+\frac{j^2}{k}<2j.
\end{equation}
Consequently $m\leq j-1$.  Substituting into
\eqref{nr:length-specialized} yields
\begin{equation}\label{nr:length-parity}
 0\leq N=3(1-j)+2m\leq1-j.
\end{equation}
It follows that $j=1$, then $m=0$ and $N=0$, so $r=s=k+1$.
This integral bound excludes every $j\geq2$, including the boundary
values for $k=3$ and $k=4$, without a separate numerical enumeration.

In every case $r=s=k+1$ and $N=0$.  Since $n_1,n_2\geq0$, they
both vanish.  Finally
\[
 2m_T=(dK-H)^2-K(dK-H)
      =d(d-1)k-2(d-1)(k+1),
\]
which equals $0$ for $(k,d)=(1,4)$ and $2(k-2)$ for $d=3$.
This proves all assertions.
\end{proof}

For reference, the possibilities proved above are
\begin{equation}\label{nr:numerical-table}
\begin{array}{c|cc|cc|c|cc}
 k&d&p&KH&H^2&n_1+n_2&m_T&m_H\\\hline
 1&4&0&2&2&0&0&0\\
 2&3&0&3&3&0&0&0\\
 3&3&1&4&4&0&1&0\\
 4&3&2&5&5&0&2&0\\
 5&3&3&6&6&0&3&0
\end{array}.
\end{equation}

\subsection{An effective pair and the sign of its contribution}

\begin{lemma}\label{nr:effective-pair}
The actual line bundles
\[
 H=K+h',\qquad T=(d-1)K-h'
\]
are effective and satisfy \eqref{nr:survivors}, with $\alpha H\ne0$.
\end{lemma}

\begin{proof}
Equation~\eqref{nr:h-prime} gives
$KH=H^2=k+1$.  The calculations in the preceding proof then give
$m_H=0$, $m_T=p$, and $N=0$.  Consequently
\[
 \chi(H)=1,\qquad \chi(T)=p+1>0.
\]
Let $B_0=\pi^*K_0$, a nef divisor on $X$.  Its pairings with the
Serre-dual line bundles are
\[
 B_0(K-H)=-a<0,\qquad
 B_0(K-T)=(2-d)k_0+a<0.
\]
For the second inequality, in the odd case its value is $-1$ when
$k_0=1$ and $1-k_0<0$ when $k_0\geq2$; in the even case it is
$-8+2=-6$.  Thus neither $K-H$ nor $K-T$ has a section, since an
effective divisor has nonnegative intersection with a nef divisor.
Serre duality gives $H^2(H)=H^2(T)=0$, and Riemann--Roch now yields
\[
 h^0(H)=\chi(H)+h^1(H)>0,\qquad
 h^0(T)=\chi(T)+h^1(T)>0.
\]
Both line bundles are therefore effective.  Finally
$\alpha H=\alpha h'=\alpha^2<0$ by
\eqref{nr:alpha-identities}.
Thus this pair occurs in the fixed locus with both point subschemes
empty, by Lemma~\ref{lem:rank-one-family}.
\end{proof}

\begin{proof}[Proof of Proposition~\ref{prop:nonrank}]
Use the data \eqref{nr:data} and pair the Chow relation of
Theorem~\ref{thm:chow-relation} with $\alpha\boxtimes\alpha$ in
cohomology.  We compute the contributions using
\eqref{rel:fiber-contribution} and \eqref{rel:pairing}.

For any contribution with $\alpha H=0$, the class $\alpha$ is
orthogonal to $K,L_1,L_2,T,H$.  Thus its external-product part has
zero pairing, and it contributes $c_*\alpha^2$ by
Proposition~\ref{prop:universal}.

For every remaining contribution, Lemma~\ref{nr:surviving-numerics}
gives $n_1=n_2=0$.  Its chosen diagonal coefficient $c_*$ is therefore
zero by Proposition~\ref{prop:universal}.  Its point Hilbert factors
and fixed Euler factors are trivial.  Formula~\eqref{eq:insertion}
gives
\[
 \mu_\alpha=(\alpha H)(-w_1+w_2+x_1-x_2)
             =(\alpha H)(-1+x_1-x_2).
\]
Since $m_T=p$ and $m_H=0$, its full pairing is
\begin{align*}
 &[x_1^p x_2^0]\,
 \mathcal R(x_1,x_2)(x_1+x_2)^p
           (\alpha H)^2(-1+x_1-x_2)^2\\
 &\hspace{35mm}=(\alpha H)^2\mathcal R(0,0).
\end{align*}
Here all denominators in $\mathcal R$ have nonzero constant term.
Setting $x_2=0$ and extracting $x_1^p$ therefore uses only the constant
term of the factors multiplying $x_1^p$.

We compute this constant exactly.  Since there are no point factors,
\eqref{eq:moving-ratio} reduces at $(x_1,x_2)=(0,0)$ to the scalar
weight powers.  The $i=2,j=1$ weight is $1$, and the $i=1,j=2$
weight is $-1$, with
\[
 P_{12}=T+K,\qquad B_{12}=T-H+K.
\]
Thus
\begin{align*}
 \mathcal R(0,0)
 &=(-1)^{\chi(T-H+K)-\chi(T+K)}\\
 &=(-1)^{-TH+(H^2-KH)/2}
  =(-1)^{-TH}.
\end{align*}
The second equality follows by expanding Riemann--Roch, and the last
uses $H^2=KH$.  Moreover
\[
 TH=(dK-H)H=(d-1)(k+1)
 =\begin{cases}6,&k=1,\\2(k+1),&2\leq k\leq5,\end{cases}
\]
which is even.  Hence every such term has
\begin{equation}\label{nr:moving-constant}
 \mathcal R(0,0)=1.
\end{equation}

Let the sum below run over the effective ordered line decompositions
occurring in the fixed locus with $\alpha H\ne0$, retaining distinct
torsion twists as distinct terms.  This is a finite sum, as the fixed
locus has only finitely many component data.  It is nonempty by
Lemma~\ref{nr:effective-pair}.  Summing the two kinds of contributions
therefore gives
\begin{equation}\label{nr:positive-coefficient}
 0=c\alpha^2+\sum_{\alpha H\ne0}(\alpha H)^2,
 \qquad
 c=-\frac{\displaystyle\sum_{\alpha H\ne0}(\alpha H)^2}{\alpha^2}>0.
\end{equation}
In this equality $c$ is the coefficient already chosen in the Chow
relation: the second group of terms has $c_*=0$, and the first group
accounts for its entire sum.  The strict positivity follows from
$\alpha^2<0$ and the nonempty sum of positive real squares.

Corollary~\ref{rel:correspondence} now gives
the rational degree isomorphism on $X$.  Lemma~\ref{lem:birational}
passes it to $X_0$, proving the Proposition.
\end{proof}
